\section{Threats to Validity and Limitations}
\label{sec:threats}

\noindent\textbf{External validity.}
Our claims are scoped to the task and corpus defined in \autoref{sec:task-methodology}: 95 HeCBench programs exposing 254 first-invocation kernels, four NVIDIA GPUs, and per-precision DRAM-level \rai on a kernel's first invocation.
Our corpus is a meaningful but selective sample of CUDA and OpenMP software, and the architectural generalization claims extend only to these device families.
The dataset does not cover AMD or Intel GPUs, other programming models such as HIP and SYCL, steady-state warm-cache behavior, or end-to-end application performance, and the results should not be extrapolated beyond execution-free first-invocation triage.

\noindent\textbf{Construct validity.}
Ground truth comes from profiler counters rather than source-level reasoning, so the byte denominator is DRAM-visible traffic rather than all source-level loads and stores; cache write-back behavior, replay effects, and profiler noise can therefore influence the target quantity, which we mitigate but do not eliminate by averaging three first-invocation profiles.
Closed-model behavior can also drift across snapshots, and a single query per input cannot establish reliability.
We address the latter directly: rather than disclaim run-to-run variability, RQ3 measures it on a stratified repeated-query panel and reports where single-shot use is and is not safe.

\noindent\textbf{Baseline limitations.}
The static baselines are deliberately modest.
The deterministic references use approximate lexical and instruction-mix counts rather than compiler-grade analyzers, and the learned tree baselines are evaluated with grouped cross-validation on this corpus rather than a separately curated external split (though they remain stable across five seeds).
Stronger compiler-driven or larger-scale learned baselines could therefore outperform the references reported here.

\noindent\textbf{Failure-analysis validity and dataset leakage.}
The feature/error analysis is exploratory: its source-feature labels come from an auxiliary LLM-voting pipeline rather than human annotation, so we treat the feature/error association as descriptive evidence about recurring error patterns, not statistically validated causal evidence.
Separately, because the models are off-the-shelf and trained on large web corpora, some kernels or their close variants may have appeared in pretraining data, and we cannot rule out such contamination completely.
Although we doubt any of the sampled LLMs were exposed to the profiling data and task in our corpus, an additional study of kernel variants with source perturbations would clarify the extent to which the LLMs generalize to more kernels.
Both are reasons to frame this work as an empirical study rather than a claim about intrinsic reasoning ability in isolation.
